%%
% 引言或背景
% 引言是论文正文的开端,应包括毕业论文选题的背景、目的和意义;对国内外研究现状和相关领域中已有的研究成果的简要评述;介绍本项研究工作研究设想、研究方法或实验设计、理论依据或实验基础;涉及范围和预期结果等。要求言简意赅,注意不要与摘要雷同或成为摘要的注解。
%%

\firstchapter{绪论}
\raggedbottom

\section{研究背景与意义}
随着足式机器人相关技术不断发展，这类系统在复杂环境中的运动能力已经有了明显提升。相比四足结构，双足机器人在结构和控制上都会更敏感一些，比如高重心、多自由度以及强动力学耦合等特点，会让系统在行走、转向以及受到外界扰动时更容易失去平衡。以 Unitree G1 这一类双足机器人为例，其支撑多边形本身就比较狭窄，而且在运动过程中还会不断变化，这就使得系统对建模误差、外界干扰以及传感器噪声都比较敏感。在这种情况下，想要实现稳定控制，其难度会明显增加，尤其是在动态运动场景中，这种问题会被进一步放大。

从已有研究来看，传统的双足机器人控制方法大多是基于简化动力学模型，比如零矩点（Zero Moment Point, ZMP）理论和线性倒立摆模型（Linear Inverted Pendulum Model, LIPM）\cite{Kajita2003}，再结合模型预测控制（Model Predictive Control, MPC）\cite{Kuindersma2016}或者全身控制（Whole-body Control, WBC）\cite{Sentis2010}来实现轨迹跟踪和姿态调节。这一类方法在规则环境中往往表现较为稳定，但一旦环境变复杂，比如出现非理想接触、复杂地形或者多自由度强耦合的情况，就容易暴露出一些问题，比如建模困难、计算开销大，以及对扰动的适应能力不够强。随着机器人自由度不断增加（例如 Unitree G1 已经达到 29 自由度），控制问题本质上会演变成一个高维非线性优化问题，这不仅会让求解变得更加困难，也很难满足实时控制的要求，从工程应用的角度来看，这一点是一个比较明显的瓶颈。

近年来，深度强化学习（Deep Reinforcement Learning, DRL）逐渐被引入到机器人运动控制领域，提供了一种不同于传统方法的思路。相比依赖精确动力学模型的控制方法，强化学习更强调通过与环境交互来学习策略，因此在高维、强非线性以及不确定性较强的场景下，往往能够表现出更好的适应能力。同时，训练完成后的策略在执行时只需要进行前向计算，这在计算效率上也具有一定优势。不过，把深度强化学习直接应用到高自由度双足机器人上，也并不是一件很“直接”的事情。一方面，动作空间维度较高，会明显增加策略学习的难度；另一方面，全身自由度带来的惯性耦合效应，也可能对步态稳定性产生更加复杂的影响。因此，如何在保证系统物理一致性的前提下，一方面降低控制维度，另一方面又尽量保留全身动力学特性，从而提升策略的鲁棒性，就成为一个需要重点考虑的问题。

基于这样的背景，本文以高自由度双足机器人为研究对象，尝试在全身动力学约束条件下，探索更加稳定的步态控制策略学习方法。具体来说，一方面通过构建全身动力学模型来保留系统的耦合特性，另一方面通过控制降维的方式降低学习复杂度，希望在两者之间找到一个相对平衡的方案，从而提升机器人在复杂环境中的动态稳定性和抗扰能力。相关研究在一定程度上也可以为后续的仿真到现实迁移（Sim-to-Real）提供参考。

\section{国内外相关技术研究及演进现状}
近年来，基于深度强化学习的机器人运动控制方法发展较快，相关研究最早主要集中在四足机器人平台上，并逐步向更复杂的双足人形系统拓展。从已有工作来看，四足机器人由于结构相对稳定，通常更适合作为强化学习方法的验证平台。在这一方向上，Hwangbo等人\cite{Hwangbo2019}较早将端到端强化学习方法应用于实际机器人控制，通过构建包含接触信息与执行器特性的状态空间，并结合执行器动力学建模和高频奖励反馈机制，实现了复杂地形下的稳定行走。在此基础上，域随机化（Domain Randomization）方法\cite{Peng2018}以及观测历史编码等策略被逐步引入训练过程，这些方法在一定程度上提升了策略对传感器噪声、动力学偏差及环境不确定性的适应能力，同时也增强了从仿真到真实环境（Sim-to-Real）的迁移能力\cite{Kumar2021}。

国内在该领域同样开展了较多研究。例如，陈学超等人\cite{Chen2023}针对双足机器人在动态行走及复杂地形适应中的控制问题，从动力学角度进行了系统分析，并指出随着关节自由度的增加，控制问题会逐渐转化为高维非线性优化问题，这也成为制约系统性能的重要因素之一。

进一步对比可以发现，相较于四足系统，双足人形机器人在强化学习控制中面临更加明显的维度挑战。一方面，高重心以及不断变化的支撑区域使系统稳定性更难维持；另一方面，多自由度带来的动作空间膨胀也显著增加了策略学习的难度。因此，早期研究通常不会直接在完整自由度空间中进行策略学习，而是倾向于采用降维或分层控制的方式来降低问题复杂度。例如，Siekmann等人\cite{Siekmann2021}提出了分层强化学习框架，将高层的步幅与步频决策与底层关节控制进行解耦，从而在一定程度上提高了训练稳定性与收敛效率；Castillo等人\cite{Castillo2021}则通过强化学习构建鲁棒反馈控制方法，验证了该类方法在抗干扰能力方面的潜力。同时，邱建斌等人\cite{10}的综述研究指出，尽管分层结构或降维方法能够降低学习难度，但在复杂转向、摆臂补偿以及姿态自平衡等任务中，这类方法在一定程度上限制了机器人全身协同能力的发挥。为了进一步提升控制效果，近年来有研究开始引入动作先验来辅助策略学习。例如，Peng等人\cite{Peng2018}通过利用动物运动视频和运动捕捉数据作为先验，引导机器人学习更加自然且高效的步态。这类方法能够加快训练过程，但同时也会在一定程度上引入对先验数据的依赖。

近两年，随着计算集群算力和最新训练范式的突破，越来越多学者开始挑战无任何参考运动先验的全身强化学习控制方法。其中，Radosavovic等\cite{Radosavovic2023}实现了高复杂双足人形机器人的端到端控制，并展示了其在扰动环境下的恢复能力；Duan等\cite{Duan2024}则探索了基于多模态感知的全身强化学习控制架构，体现了无先验条件下的策略学习潜力。

综上分析，虽然通过减少自由度或引入结构化先验可以加速收敛过程，但由于Unitree G1系统具有29个自由度及强耦合动力学特性，传统方法已难以充分发挥全身协调能力。因此，本文采用全自由度耦合建模与动作空间解耦结合的方式，在降低学习复杂度的同时保留系统动态特性，以提升复杂场景下的适应能力。  

\section{本文主要研究内容与章节安排}
本文以 Unitree G1 双足人形机器人为研究对象，主要围绕高自由度系统下的稳定步态控制问题展开研究。从整体思路来看，一方面需要考虑全身动力学耦合带来的影响，另一方面又要控制策略学习的复杂度，这两点之间需要做一定的权衡。在具体实现上，首先在 NVIDIA Isaac Gym 仿真环境中构建了包含 29 自由度的全身动力学模型，并将双足机器人的运动控制过程建模为强化学习中的马尔可夫决策过程。在此基础上，采用近端策略优化算法（PPO）进行策略训练，使模型能够通过与环境交互逐步学习稳定的控制策略。针对高维动作空间带来的训练难度问题，本文没有直接对全部自由度进行控制，而是引入了控制降维的设计思路。具体来说，策略网络只输出下肢关键关节的控制信号，而上半身自由度则通过底层 PD 控制器进行约束。这样一来，在保留全身动力学耦合特性的同时，也在一定程度上降低了策略学习的复杂度。在奖励函数设计方面，本文从多个角度对策略进行引导，包括速度跟踪、姿态稳定、足端运动约束以及动作平滑性等因素，通过多目标组合的方式，使策略逐步学习到更加稳定且自然的步态行为。同时，在训练过程中引入扰动和随机化机制，以增强策略在复杂环境和非理想条件下的适应能力。在实验部分，基于大规模并行仿真环境，对不同自由度建模方式下的控制策略进行了对比分析，重点从收敛特性、运动稳定性以及抗扰能力等方面进行评估，从而验证所提出方法的有效性。

在上述研究基础上，本文的整体结构安排如下：第 1 章为绪论，对研究背景及相关工作进行介绍；第 2 章主要阐述强化学习在连续控制问题中的理论基础；第 3 章介绍系统建模方法、策略网络结构以及奖励函数设计；第 4 章给出仿真实验设置及结果分析；第 5 章对全文工作进行总结，并对后续研究方向进行展望。
